\Needspace{31\baselineskip}\vspace{6mm}
\pdfbookmark[1]{Letter from Gauss to Schumacher. 1836}{letter546}
\begin{center}
{\large\itshape Letter from Gauss to Schumacher.\par}
\vspace{5mm}\hrule\vspace{2mm}
{\small \textsc{Schumacher}. Astronomische Nachrichten No.\ 310. 11 June 1836.\par}
\vspace{2mm}\hrule\vspace{5mm}
\end{center}
\begin{flushright}Göttingen. 23 April 1836.\end{flushright}
The magnetic needle made extraordinary movements this morning, greater even than those on 7 February 1835. This prompted me to take a few sets of readings at the astronomical observatory while Dr \textsc{Goldschmidt} recorded them at the magnetic \SourceLeaf{547}observatory. The uniform course of the movement was confirmed so beautifully here that I ventured to deduce the difference between the two clock readings from it.

\Needspace{9\baselineskip}
\begin{center}\small
\begin{tabular}{@{}p{0.82\linewidth}r@{}}
From a rapid rise, Campa was found to be ahead of Shelton by & $4'\,41''1$\\
From a fall a few minutes later & $4\quad42.4$\\\cline{2-2}
\multicolumn{1}{r}{Mean} & $4'\,41''7$
\end{tabular}
\end{center}
A direct comparison of the clocks gave
\begin{center}\small
\begin{tabular}{@{}lll@{}}
1) by a signal at the window & . . . . . . . . . . & $4'\,41''5$\\
2) by an induction pulse & . . . . . . . . . . & $4\quad41.5$
\end{tabular}
\end{center}
A splendid confirmation, then, of what I said in Astronomische Nachrichten No.\ 276 [p.\ 539 of this volume].

\clearpage
